% !TEX program = pdflatex
\documentclass[12pt,a4paper]{letter}

\usepackage[T1]{fontenc}
\usepackage[margin=1in]{geometry}
\usepackage{hyperref}
\usepackage{setspace}
\usepackage{parskip}

\hypersetup{
  colorlinks=true,
  urlcolor=blue
}

\onehalfspacing

\signature{Professor {[FULL NAME]}\\
Professor of Computer Science\\
{[Department / Faculty]}\\
{[University Name]}, {[Country]}}

\address{{[Department of Computer Science]}\\
{[University Name]}\\
{[City, State/Province]}\\
{[Country]}}

\date{27 February 2026}

\begin{document}

\begin{letter}{UK Visas \& Immigration\\
Global Talent Visa --- Tech Nation (Digital Technology)\\
United Kingdom}

\opening{Dear Sir/Madam,}

This letter serves as my endorsement of Odeyemi Olusegun Israel for the UK Global Talent Visa in Digital Technology. I should state at the outset that I do not write these letters often, and when I do, it is because the individual's work has given me genuine cause to pay attention.

I am a Professor of Computer Science at {[University Name]}, where my research interests include {[relevant areas, e.g., distributed systems, machine learning, software engineering]}. I have known Olusegun as a mentor over a period of {[X]} years and have had the opportunity to observe both his technical growth and, more recently, the research outputs he has produced independently.

\textbf{On the engineering work.} Olusegun has built a platform called AMTTP---the Anti-Money Laundering Transaction Trust Protocol---and it is, frankly, more ambitious than most projects I see coming out of funded research groups. The system spans four architectural layers: a client-facing SDK layer (TypeScript and Python, both open-sourced), a compliance orchestration engine, an offline ML training pipeline, and a blockchain smart contract layer with 18 contracts deployed on Ethereum Sepolia. The whole thing runs across 17 containerised microservices coordinated through a central gateway.

To anyone outside the field, that list might read as a collection of buzzwords. It is not. Each of those components has to actually work together, and making them do so---reliably, at scale, with deterministic compliance guarantees---is a hard systems problem. I know this because I have watched postdoctoral researchers struggle with far simpler integration challenges. Olusegun has done it alone.

The core technical insight behind AMTTP is worth explaining. In traditional finance, a suspicious transaction can be halted mid-flight. On a blockchain, once a transaction reaches finality, it is irreversible. Every existing AML tool on the market---Chainalysis, Elliptic, and their competitors---operates after this point. They flag; they do not prevent. AMTTP operates before settlement, combining ML risk scores, graph topology analysis, sanctions list screening, and configurable policy rules into a single pass/fail decision. That is a meaningful architectural contribution to the field and, as far as I am aware from surveying the literature, no other system does this at the protocol level.

\textbf{On the machine learning contributions.} Two of Olusegun's papers have caught my attention for different reasons.

The first, on \textit{Blind Spot Decomposition Theory} (BSDT), asks a question that most ML practitioners avoid because it is uncomfortable: why do well-performing models miss the fraudulent transactions they miss? His answer is structural. He decomposes missed fraud into four components---Camouflage, Feature Gap, Activity Anomaly, and Temporal Novelty---and provides evidence that these four components, taken together, account for over 95\% of the explainable variance in the missed-fraud indicator. The correction operators he derives from this decomposition recover recall by as much as 84.4 percentage points. He validates this across four different model architectures and six datasets drawn from five domains: Bitcoin transactions, Ethereum transactions, credit card fraud, NASA shuttle telemetry, mammography, and digit recognition. The work requires no retraining and no access to model internals, which makes it practically useful in a way that many theoretical contributions are not.

I want to be precise about what makes this good. Plenty of researchers can improve a metric on a benchmark. Fewer can articulate \textit{why} a model fails in a way that generalises across domains. Olusegun has done the latter, and the experimental rigour---four-way data splits to prevent label leakage, systematic comparison of 13 alternative combination strategies, false-discovery analysis---is thorough.

The second paper, on the \textit{Universal Deviation Law} (UDL), tackles anomaly detection from a different angle entirely. Rather than asking why detectors fail, it proposes a tensor-based mathematical framework for measuring deviation across multiple spectral representations simultaneously. The practical consequence is a detection method that catches anomalies which single-spectrum approaches miss by construction. It is domain-agnostic, meaning it applies equally to financial transactions, sensor data, or medical imaging.

Taken together, these two papers constitute a coherent research programme: BSDT diagnoses failure; UDL provides a more robust detection foundation. That kind of conceptual coherence rarely emerges by accident.

\textbf{On the broader system design.} Several details in Olusegun's work reflect a maturity that goes beyond raw technical ability. He has implemented a zero-knowledge proof framework (zkNAF) that allows institutions to verify KYC credentials, risk score ranges, and sanctions non-membership without exposing the underlying personal data. This is not trivial cryptographic engineering, and it addresses a real tension between GDPR privacy obligations and AML transparency requirements that regulators have not yet fully resolved. He has also integrated multi-oracle threshold signatures and replay protection---the sort of infrastructure-level security work that is essential but rarely glamorous enough to feature in papers.

The ML training pipeline itself deserves mention. It uses a Composite Teacher architecture: Autoencoder-Enhanced XGBoost (weighted at 0.4), seven FATF-aligned AML pattern detectors (0.3), and graph structural properties (0.3) to generate pseudo-labels across 2.64 million transactions. A Student model then learns from these labels. This teacher-student distillation approach is well-suited to a domain where ground-truth labels are scarce and expensive, and the weighting reflects a considered judgment about how much to trust each signal source.

\textbf{My assessment.} I have supervised and examined a good number of PhD theses and postdoctoral projects over my career. If I am honest, some of them contained less original thinking than what Olusegun has produced as an independent researcher working outside the university system. That is not a comfortable observation for an academic to make, but it is an accurate one.

His work sits at the intersection of machine learning, distributed systems, blockchain architecture, and financial regulation---each of which is a deep field in its own right. The fact that he has produced functioning software, original theory, and experimentally validated results across all four areas tells me something about both his ability and his work ethic.

The UK's financial technology sector is, by most measures, the strongest in Europe. The FCA's forthcoming regulatory framework for crypto-assets, expected by late 2027, will create significant demand for exactly the kind of compliance infrastructure Olusegun has already built. I believe he would make a substantive contribution to the UK's digital economy, and I endorse his application for the Global Talent Visa accordingly.

\closing{Yours faithfully,}

\end{letter}

\end{document}
